\section{Getting to Know Your Equipment}
\label{lab_equipment}

%\makelabheader %(Space for student name, etc., defined in master.tex)

\bigskip

\textit{Note: good preparation for this lab is looking over the introduction and chapter 2 of Paynter.  This will help you to identify what's what in your kits.}

\begin{enumerate}[wide]

\item Identify the values of some of the resistors in your kit, using the resistor color band code 
%on page 54 of Paynter.  
in Appendix \ref{resistor_code}.
What is the ``tolerance'' of most of your resistors, and what does that mean?   

\item Use your digital multimeter (DMM) to measure the resistances of a 1 k$\Omega$ and 1 M$\Omega$ resistor from your kit.  How close are they to their stated values?

\item What is the accuracy of the DMM for the resistance for the readings above?  (Look in the manual for the DMM to find this out.)  In general, which is more accurate: the color bands on the resistors, or a measurement from your DMM?

\item Measure the resistance of yourself.  Does it make a difference if you lick your fingers?  

\item What does the ``range'' button do on your DMM?  What is ``auto ranging'' and how do you get it to do that?

\item What does the min/max/avg button on the DMM do?

\item What does the yellow button on the DMM do, both in general, and specifically when you're measuring resistances?  (What does the ``
\includegraphics[height=0.2in]{equipment/sound.eps}
'' symbol mean, and when might that be useful?) \label{part_audible_dmm}

\item Use what you just found in part \ref{part_audible_dmm} to figure out which holes on your ``breadboard'' are connected to which other holes.   

\item What is the resistance of connections you make using jumper wires plugged into your breadboard.  If you can't measure the resistance precisely, can you at least set an upper limit?

\item The black banana jack on your proto boards is labeled with a triangle, which typically means ``signal common.'' Often, this is also connected to ``ground,'' that is, the Earth.  Is this actually connected to the Earth?  How can you test it for sure?  What else on your proto boards is connected to ``signal common?''

\end{enumerate}
\begin{enumerate}[wide,resume]


\begin{minipage}{.75\textwidth}

\item A ``single pole double throw'' (SPDT) switch has three terminals, as shown in the diagram to the right.  The switches S9 and S10 on your proto boards each have eight little holes for wire connections; the leftmost two are always connected to each other, as are the rightmost two and the four in the center.  Which of these sets of holes correspond to each of the three terminals?  Which holes are connected to each other when the switch is up, and when it is down?  Make a good diagram in your lab notebook; you'll need to know how to use these later! \label{part_switches}
\end{minipage}
\begin{minipage}{.24\textwidth}
%\center{\includegraphics[width=0.9\textwidth]{equipment/spdt.eps}}
\center{\includegraphics{equipment/spdt2.eps}}
\center{\includegraphics[width=0.7\textwidth]{equipment/spdt_contacts.eps}}
\end{minipage}

\item Use your DMM to measure the resistance of the switch when it is closed.  If you can't make an accurate measurement, can you at least set an upper limit on the resistance? 

\item Measure the voltage of your 5~V power supply on your proto boards.  (This means measuring the voltage between the 5 V supply and ground.)  How close to 5.0000000 volts is it?  Which is more accurate: the ``5 V'' printed on the proto board, or a measurement by your DMM?

\item Use your DMM to measure the actual range of the variable supplies $+V$ and $-V$ on your proto boards.  Are their ranges exactly 1.3 to 15 volts when you turn the knobs?

\item Set your function generator to a 1.0 kHz sine wave, and use your DMM to measure the frequency and amplitude of the oscillating voltage with the amplitude slider all the way up.   As with your 5~V supply, you will be measuring the voltage between the function generator's output and ground.  Will that voltage be AC or DC?  (Note that the function generator has a block of eight holes for the output, similar to the switches S9 and S10.  In this case, the rightmost six are all connected together; that's the output you want.  The leftmost two are the ``TTL'' signal; you'll learn about that in Lab \ref{lab_oscilloscope}.)

\item Just for fun, let's hook up the function generator to your speaker.  The speaker has two terminals, one connected to the top four holes in the block of eight, and one connected to the bottom four holes.  One speaker terminal should connect to the function generator's output.  Where should the other terminal connect to?  Once you have it hooked up, what's the highest frequency you can hear?  How does that compare with the highest audible frequency for your teacher, who spent way too much time playing in really loud, bad rock bands in college?
 
\end{enumerate}




